\documentclass[12pt,a4paper]{article}

\usepackage{fontawesome5}
\providecommand{\faPencilAlt}{\faIcon{pencil-alt}}

\usepackage{styles/couleurs}
\usepackage{styles/mise_en_page}
\usepackage{styles/boites}
\usepackage{styles/maths}
\usepackage{styles/figures}

\usepackage{enumitem}
\usepackage{array}
\usepackage{tabularx}
\usepackage{amsmath}
\usepackage{amssymb}
\usepackage{tikz}
\usepackage{xparse}
\usepackage{graphicx}
\usepackage{pdflscape}
\usepackage{colortbl}

\usetikzlibrary{calc,arrows.meta,positioning,patterns}

\setlist[itemize]{leftmargin=1.4em,itemsep=0.15em,topsep=0.2em}
\setlist[enumerate]{leftmargin=*,itemsep=0.22em,topsep=0.25em}

\newcommand{\lignepointillee}{%
  \noindent\dotfill\par
  \vspace{0.45em}
}

\newcommand{\lignesreponse}[1]{%
  \begin{tcolorbox}[
    enhanced,
    colback=white,
    colframe=NeutreC!40,
    boxrule=0.5pt,
    arc=1.5mm,
    left=3mm,
    right=3mm,
    top=2mm,
    bottom=1mm,
    before skip=2pt,
    after skip=2pt
  ]
    \foreach \i in {1,...,#1}{\lignepointillee}
  \end{tcolorbox}
}

\newcommand{\imageactivite}[3]{%
  \IfFileExists{#3}{%
    \includegraphics[width=#1,height=#2,keepaspectratio]{#3}%
  }{%
    \fbox{%
      \parbox[c][#2][c]{#1}{%
        \centering\sffamily\small
        Image à intégrer\par
        \vspace{0.3em}
        \texttt{\detokenize{#3}}
      }%
    }%
  }%
}

\newcommand{\titreseq}[1]{%
\begin{center}
  {\sffamily\LARGE\bfseries\color{TitreC}
  Séquence 1 -- Statistiques}\\[0.2em]
  {\sffamily\large\bfseries\color{SavoirC}#1}
\end{center}
}

\begin{document}

\titreseq{Cours -- Lire, résumer et représenter des données}

\section{Vocabulaire}

\definition{
\begin{itemize}
  \item La \textbf{population} est l'ensemble étudié.
  \item Le \textbf{caractère} est la question ou la propriété étudiée.
  \item Les \textbf{valeurs} ou \textbf{modalités} sont les réponses possibles.
  \item L'\textbf{effectif} d'une valeur est le nombre de fois où elle apparaît.
  \item L'\textbf{effectif total} est le nombre total de données.
\end{itemize}}

\exemple{On étudie le mode de transport de \(24\) élèves.

La population est l'ensemble des \(24\) élèves.
Le caractère étudié est le mode de transport.
Les modalités peuvent être : bus, voiture, vélo, à pied.}

\section{Fréquences}

\definition{La fréquence d'une valeur est :

\[
\boxed{
\text{fréquence}
=
\text{effectif de la valeur}
\div
\text{effectif total}
}
\]}

\propriete{Une fréquence peut s'écrire sous forme de fraction,
de nombre décimal ou de pourcentage.

La somme des fréquences d'une répartition complète vaut \(1\),
soit \(100\%\).}

\exemple{Dans une classe de \(25\) élèves, \(10\) viennent à pied.

\[
10\div25=0{,}4=40\%.
\]}

\section{Moyenne}

\definition{La moyenne d'une série est :

\[
\boxed{
\text{moyenne}
=
\text{somme des valeurs}
\div
\text{nombre de valeurs}
}
\]}

\exemple{Pour la série \(8\), \(10\), \(12\), \(14\) :

\[
(8+10+12+14)\div4=11.
\]}

\remarque{La moyenne est le nombre qui pourrait remplacer toutes les valeurs
si elles devenaient égales tout en conservant la même somme totale.

Elle ne correspond pas nécessairement à une valeur réellement observée.}

\subsection{Moyenne à partir d'un tableau d'effectifs}

\methode{On multiplie chaque valeur par son effectif,
on additionne les produits, puis on divise par l'effectif total :

\[
\boxed{
\text{moyenne}
=
\text{somme des produits valeur}\times\text{effectif}
\div
\text{effectif total}
}
\]}

\exemple{
\[
\begin{array}{c|ccc}
\text{Valeur} & 5 & 10 & 15\\
\hline
\text{Effectif} & 2 & 3 & 1
\end{array}
\]

\[
(5\times2+10\times3+15\times1)\div6
=
55\div6
\approx9{,}2.
\]}

\section{Médiane}

\definition{Une médiane partage une série ordonnée en deux groupes
de même effectif :

\begin{itemize}
  \item au moins la moitié des valeurs sont inférieures ou égales à la médiane ;
  \item au moins la moitié des valeurs sont supérieures ou égales à la médiane.
\end{itemize}}

\methode{\textbf{Étape 1 :} ranger les valeurs dans l'ordre croissant.

\textbf{Étape 2 :} repérer la ou les valeurs centrales.

\begin{itemize}
  \item Si l'effectif est impair, la médiane est la valeur centrale.
  \item Si l'effectif est pair, on prend la moyenne des deux valeurs centrales.
\end{itemize}}

\exemple{\textbf{Effectif impair}

\[
4\quad7\quad9\quad10\quad12.
\]

La valeur centrale est \(9\), donc la médiane est \(9\).}

\exemple{\textbf{Effectif pair}

\[
3\quad5\quad6\quad8\quad11\quad14.
\]

Les deux valeurs centrales sont \(6\) et \(8\).

\[
(6+8)\div2=7.
\]

La médiane est \(7\).}

\remarque{La médiane peut être déterminée à partir de données brutes,
d'un tableau d'effectifs ou d'un diagramme en bâtons,
à condition de repérer les positions dans la série ordonnée.}

\section{Moyenne ou médiane ?}

\propriete{
\begin{itemize}
  \item La moyenne utilise toutes les valeurs et conserve la somme totale.
  \item La médiane dépend de la position des valeurs dans la série ordonnée.
  \item Une valeur extrême peut modifier fortement la moyenne.
  \item La médiane est généralement moins sensible aux valeurs extrêmes.
\end{itemize}}

\exemple{Dans la série :

\[
1500\quad1550\quad1600\quad1650\quad1700\quad9000,
\]

le salaire de \(9000\) euros augmente fortement la moyenne,
mais il ne modifie pas les deux valeurs centrales.}

\remarque{Le choix de l'indicateur dépend de la question posée.

\begin{itemize}
  \item Pour une répartition équitable d'un total, la moyenne est souvent pertinente.
  \item Pour décrire une valeur centrale malgré une valeur extrême,
        la médiane peut être plus informative.
\end{itemize}}

\section*{Checklist}

\begin{tcolorbox}[
  enhanced,
  colback=SavoirBg,
  colframe=SavoirC!65,
  boxrule=0.7pt,
  arc=2mm
]
Avant de terminer, je vérifie :

\begin{itemize}
  \item[\(\square\)] que mes effectifs ont le bon total ;
  \item[\(\square\)] que mes fréquences totalisent \(100\%\) ;
  \item[\(\square\)] que ma série est rangée avant de chercher la médiane ;
  \item[\(\square\)] que j'ai pris les bonnes positions centrales ;
  \item[\(\square\)] que la moyenne possède la bonne unité ;
  \item[\(\square\)] que j'ai interprété l'indicateur dans le contexte ;
  \item[\(\square\)] que mon graphique comporte titre, légende et informations utiles.
\end{itemize}
\end{tcolorbox}

\newpage

\section{Tableaux et représentations}

\remarque{Une représentation est choisie pour répondre à une question.

\begin{itemize}
  \item Un tableau organise précisément les données.
  \item Un diagramme en bâtons compare des effectifs ou des fréquences.
  \item Un diagramme circulaire montre la part de chaque catégorie
        dans un total.
\end{itemize}}

\subsection{Diagramme circulaire}

\propriete{Dans un diagramme circulaire, les angles sont proportionnels
aux effectifs ou aux fréquences.

Le cercle entier représente :

\[
\text{effectif total}
\longleftrightarrow
100\%
\longleftrightarrow
360^\circ.
\]}

\methode{Pour calculer l'angle d'un secteur :

\[
\boxed{
\text{angle}
=
\text{effectif de la modalité}
\times360
\div
\text{effectif total}
}
\]

ou

\[
\boxed{
\text{angle}
=
\text{pourcentage}
\times360
\div100
}
\]}

\exemple{Une catégorie représente \(25\%\) du total.

\[
25\times360\div100=90.
\]

Le secteur mesure \(90^\circ\).}

\aretenir{Pour contrôler un diagramme circulaire :

\begin{itemize}
  \item la somme des pourcentages doit valoir \(100\%\) ;
  \item la somme des angles doit valoir \(360^\circ\) ;
  \item un secteur plus grand doit représenter une part plus grande ;
  \item le titre, la légende, la source et l'effectif total doivent être précisés.
\end{itemize}}



\end{document}
